\section{Results and Discussion}
\label{sec:results}

\subsection{Dataset and protocol}
The corpus contains 55 clean documents, paragraphs from the READMEs of ten
open-source projects (Node.js, axios, Redis, FastAPI, pandas, Express,
Django, NumPy, requests, Flask), and 60 poisoned documents (15 per
category): direct instruction injection, fabricated system-authority
messages, on-topic ``corpus poisoning'' with harmful advice and no trigger
phrases, and garbled text. The embedding detector's reference corpus and
its evaluation set are disjoint: 70\% of clean documents (38) form the
trusted baseline, and the remaining 17, with all 60 poisoned documents,
are scored (77 total). Provenance for every document is logged. Thresholds
and calibration bounds were chosen on these same 77 documents, so all
numbers below are development-set results and are optimistic (see
Section~\ref{sec:limitations}).

\subsection{Fusion method comparison}
Table \ref{tab:fusion} summarizes the four fusion strategies.

\begin{table}[t]
\caption{Fusion method comparison (77-document development set)}
\label{tab:fusion}
\centering
\begin{tabular}{@{}lccc@{}}
\toprule
Method & Precision & Recall & F1 \\
\midrule
Naive fixed fusion & 1.00 & 0.65 & 0.79 \\
Learned fusion (raw scores) & 0.86 & 0.90 & 0.88 \\
Calibrated fixed fusion & 0.98 & 0.97 & 0.97 \\
Learned fusion (calibrated) & 0.97 & 0.98 & 0.98 \\
Staged pipeline & 0.94 & 1.00 & 0.97 \\
\bottomrule
\end{tabular}
\end{table}

The gap between naive and calibrated fusion is explained by score range.
The embedding detector's raw scores occupied $[0.180, 0.430]$, while the
pattern and perplexity detectors each spanned the full $[0,1]$ range.
Unweighted averaging therefore compresses the embedding detector by
construction, though a per-detector threshold sweep shows it is the
strongest signal (F1 = 0.92 alone at its best threshold, against 0.55 and
0.49 for pattern and perplexity). Calibrating each detector to $[0,1]$
before averaging raises F1 from 0.79 to 0.97 with no change to any
detector's scoring logic.

The learned fusion corroborates this. A logistic regression on raw scores
reaches only F1 = 0.88 and gives the embedding detector the smallest
coefficient of the three (1.09, against 1.82 and 1.51) despite its being
the most discriminative signal. Trained on calibrated inputs the same model
reaches F1 = 0.98 and ranks the embedding detector first (coefficient
2.67). A hand-derived calibration argument and a cross-validated model thus
converge on the same conclusion.

\subsection{Failure map by attack category}
Table \ref{tab:failuremap} reports, per category, the fraction of poisoned
documents caught by each detector at its own best threshold, and by
calibrated fusion at threshold 0.20.

\begin{table}[t]
\caption{Per-category detection rate}
\label{tab:failuremap}
\centering
\begin{tabular}{@{}lcccc@{}}
\toprule
Category & Pattern & Embed. & Perplex. & Fusion \\
\midrule
Injection & 67\% & 73\% & 33\% & 100\% \\
Authority & 73\% & 93\% & 0\% & 93\% \\
Corpus poisoning & 13\% & 100\% & 0\% & 93\% \\
Garbled/perplexity & 0\% & 100\% & 100\% & 100\% \\
\bottomrule
\end{tabular}
\end{table}

No single detector covers every category, the central argument for an
ensemble. Corpus poisoning, harmful advice in ordinary technical language
with no override phrases, is caught almost entirely by the embedding
detector. As the cross-domain result below shows, part of that success is
likely topic distance from the README baseline rather than detection of
malicious intent.

\subsection{A regex gap found and closed by failure analysis}
Inspection of documents missed by calibrated fusion identified two
injection documents containing ``ignore \textit{your} previous
instructions,'' which the original regular expression did not match because
it allowed only one optional modifier before ``instructions.'' Generalizing
the pattern to accept a possessive pronoun raised injection-category
detection from 87\% to 100\% and calibrated fusion F1 from 0.95 to 0.97,
verified by re-running the full corpus to confirm no new false positives.

\subsection{Staged pipeline and latency}
Running all three detectors on every document costs 357.4\,ms on average in
our latest run. The staged pipeline (pattern, then embedding, then
perplexity only if inconclusive) costs 105.3\,ms, a 70.5\% reduction, with
F1 = 0.97. The embedding stage ran on 70\% of documents and the perplexity
stage on 17\%. The most accurate single signal need not be the most
expensive to deploy, provided cheaper signals first resolve the documents
they can. This is a replay of latencies logged in one evaluation run; an
earlier run gave 72.5\%, so the figure is indicative rather than precise.

\subsection{Adversarial red-teaming}
We wrote ten documents meant to defeat all three detectors: each avoids
the original pattern trigger phrases, is fluent to avoid a high perplexity
score, and imitates README vocabulary to keep embedding distance small.
Calibration bounds were frozen from the original evaluation rather than
recomputed, to avoid circularity. In the first round, nine of ten were
caught. The evasion, \texttt{adv\_010}, opened with ``this project is
distributed under the MIT license'' and then described build artifacts
published with default demo credentials. We call this
\textit{boilerplate anchoring}: prefixing a malicious clause with text so
common in the reference domain that it pulls the whole-document embedding
toward the baseline and dilutes the payload (combined score 0.106 against
a threshold of 0.15). We did not find prior work describing this effect in
this setting.

We then added two credential-exposure patterns (default/preconfigured
credentials, hard-coded secrets). They match three of the ten adversarial
documents and none of the 55 clean or 12 cross-domain documents, and all
ten are now caught (\texttt{adv\_010}: 0.406). Because the patterns were
written after reading \texttt{adv\_010}'s wording, this is a fix for a
found gap, \textbf{not} independent evidence of robustness; the embedding
and perplexity detectors alone still miss that document.

\subsection{Calibration fragility}
\label{sec:calibration-fragility}
Min-max bounds for the embedding detector are set by one document at each
extreme. Removing the single document that sets a bound shifts the lower
bound from 0.180 to 0.205 and the whole dataset's normalized embedding
scores by 2.1\% on average, roughly 1.8x more than under 5th/95th-percentile
bounds (1.1\%). The pattern and perplexity detectors are unaffected, since
their scores tie at 0 and 1. More stable calibration is not automatically
better detection: under quantile bounds \texttt{adv\_010} receives a lower
embedding contribution (0.215) than under min-max (0.318), so a switch
would require re-tuning the threshold.

\subsection{Cross-domain false positives}
Twelve benign documents (HR, support and office policies) were scored with
the same baseline, frozen bounds and threshold. All 12 were flagged, each
driven mainly by the embedding detector. Their median embedding score
(0.37) is essentially that of the poisoned documents (0.36) and well above
that of in-domain clean documents (0.27); only 1 of 17 held-out in-domain
clean documents was flagged at threshold 0.20. The baseline contains only
software READMEs, so any document in a different style looks anomalous even
when harmless. The embedding detector therefore partly measures distance
from the baseline domain, not malicious intent, and the in-domain F1 of
0.97 must not be read as domain-general.

\subsection{Negative result: segment-level scoring}
To address boilerplate anchoring we scored each sentence independently with
three aggregations (maximum, a blend of maximum and mean, and a high-threshold
veto). None helped: maximum aggregation raised clean false positives from 4
to 9 of 17 at the original threshold, and after re-tuning the peak F1 values
were 0.90, 0.90 and 0.92, none above the 0.92 of the unmodified
document-level embedding detector. The highest segment score of
\texttt{adv\_010} (0.319) was below the maximum segment score of 5 of 17
clean documents. Sentence-level scores are noisier than the document-level
average, so the anchoring weakness is better addressed at the pattern level
or by a different calibration scheme than by finer-grained embeddings.
